\SourceLeaf{20}

Die Umkehrung der vorstehenden Entwicklung lehrt, dass man aus den \textsc{Gauss}'schen Angaben im
Falle $q=3n+1$ zunächst auf die Proportion:
\[
\left[\dfrac{\varkappa_1}{\pi_1}\right]:\left[\dfrac{\pi_1}{\varkappa_1}\right]=\left[\dfrac{\varkappa_1}{\pi_2}\right]:\left[\dfrac{\pi_2}{\varkappa_1}\right]
\]
schliessen kann. Unter nochmaliger Benutzung des in dieser Proportion ausgesprochenen Gesetzes findet man:
\[
\dfrac{\left[\dfrac{\varkappa_1}{\pi_1}\right]}{\left[\dfrac{\pi_1}{\varkappa_1}\right]}=\dfrac{\left[\dfrac{\varkappa_2}{\pi_1}\right]}{\left[\dfrac{\pi_1}{\varkappa_2}\right]}=\dfrac{\left[\dfrac{\varkappa_1}{\pi_2}\right]^2}{\left[\dfrac{\pi_2}{\varkappa_1}\right]^2}
\quad\text{und also}\quad
\dfrac{\left[\dfrac{\varkappa_1}{\pi_2}\right]}{\left[\dfrac{\pi_2}{\varkappa_1}\right]}=\dfrac{\left[\dfrac{\varkappa_1}{\pi_2}\right]^3}{\left[\dfrac{\pi_2}{\varkappa_1}\right]^2}.
\]
Da hieraus $\left[\dfrac{\varkappa_1}{\pi_2}\right]=\left[\dfrac{\pi_2}{\varkappa_1}\right]$ hervorgeht, so ist auch für $q=3n+1$ in \textsc{Gauss}' Angaben das allgemeine Reciprocitätsgesetz in seiner einfachsten Form implicite enthalten.

Unter [4] sind diejenigen Aufzeichnungen gesammelt, durch welche \textsc{Gauss} seine cubischen Reciprocitätstheoreme bewiesen hat. Da diese Aufzeichnungen mehreren losen Zetteln ohne Datumangabe entnommen sind, so lässt sich die Entstehungszeit der fraglichen Beweise nicht mit Sicherheit angeben. Es ist
aber wahrscheinlich, dass die Entwicklungen alsbald nach Gewinnung der »Observatio« und der »Criterien«
durchgeführt sind; sie schliessen sich nämlich im Gedankengang sehr eng an die ersten Angaben unter [2]
und übrigens an den Art. 358 der Disquisitiones arithmeticae an, und andererseits ist \textsc{Gauss} zufolge seiner
eigenen Aussage\HistoricalNote{\textsuperscript{*)}}{*)}{Siehe den Art. 33 der am 24. August 1808 der Göttinger Societät vorgelegten Abhandlung »Summatio quarumdam serierum singularium«.} jedenfalls vor Mitte des Jahres 1808 im Besitze des sechsten Beweises des quadratischen
Reciprocitätsgesetzes gewesen, welcher mit den hier gemeinten Beweisen am nächsten verwandt ist.

Dem \textsc{Gauss}'schen Beweise des cubischen Reciprocitätssatzes liegt die Kreistheilung zu Grunde. Die
einzelnen Formeln unter [4] können entweder unmittelbar aus Art. 358 der Disquisitiones arithmeticae entnommen werden, oder sie lassen sich doch mit sehr geringer Mühe aus den dortigen Entwicklungen ableiten.
Die Congruenz:
\[
P^{pp}\equiv r^{pp}+Rr^{ppg}+\ldots\equiv P(R^{\mathrm{ind.}\,p})
\]
bezieht sich auf den Modul $p$ und entspringt aus dem Umstande, dass alle von 1 verschiedenen Polynomialcoefficienten der Potenz $pp$ durch $p$ theilbar sind. Die Fallunterscheidung
\[
\left\{\dfrac{M-N\sqrt{-27}}{M+N\sqrt{-27}}\right\}^{\frac13(pp-1)}\equiv 1,\quad -\tfrac12+\tfrac12\sqrt{-3},\quad -\tfrac12-\tfrac12\sqrt{-3}\qquad\;(\mathrm{mod.}\ p)
\]
correspondirt sowohl für $p=3m-1$ wie $p=3m+1$ genau der oben unter I, II und III getroffenen
Unterscheidung $B\equiv 0\;(\mathrm{mod.}\ q)$, etc. Die Zuordnung jener drei Congruenzen zu den folgenden
\[
p^{\frac13(n-1)}\equiv 1,\quad -\tfrac12-\tfrac12\dfrac{M}{3N},\quad -\tfrac12+\tfrac12\dfrac{M}{3N}\qquad\;(\mathrm{mod.}\ n),
\]
wie sie in [4] entwickelt wird, enthält demnach bereits einen vollständigen Beweis der \textsc{Gauss}'schen Angaben
und damit implicite des allgemeinen cubischen Reciprocitätssatzes.

Die weiteren Ausführungen unter [4] betreffen die Criterien über die cubischen Reste 2 und 3.
Wegen des cubischen Restes 3 bemerke man nur, dass (unter Benutzung der in [4] gebrauchten Bezeichnung)
$p^3\equiv p, p'$ oder $p''\;(\mathrm{mod.}\ 3)$ zutrifft, je nachdem ind.\,$3\equiv 0, 1$ oder $2\;(\mathrm{mod.}\ 3)$ stattfindet. Die Bedingung
$N\equiv 0\;(\mathrm{mod.}\ 3)$, damit 3 cubischer Rest von $n$ ist, ergiebt sich daraufhin unmittelbar aus den Schlussangaben von [4].

